\chapter*{Use Cases} \label{use_cases}
\noindent

\subsection*{1. Edge Neural Network Inference (Embedded Variant)}

\textbf{Context:} Neural network models deployed on edge IoT devices (microcontrollers, embedded GPUs, FPGAs) store quantized weight tensors in on-chip SRAM or off-chip DRAM/NVM.
Memory bit-flip rates increase with device aging, voltage scaling, and radiation exposure in harsh environments (automotive, aerospace, medical).

\textbf{MECC application:} The embedded metadata variant embeds CRC hash bits directly into the weight values via MILP co-optimization.
At zero storage overhead, the edge device stores a standard-sized weight file while retaining full MECC correction capability.
When a bit flip corrupts a weight, the on-device solver runs the coverage-ordered correction pass and restores the affected weights before inference proceeds.
The Feistel permutation ensures that burst errors from localized SRAM wear-out are equalized across the grid and corrected as reliably as isolated random errors.

\textbf{Key benefit:} Zero additional SRAM required; inference accuracy is preserved under BER up to 5\%.

\subsection*{2. NAND Flash Storage for Data Archiving}

\textbf{Context:} NAND flash memory degrades with write/erase cycles, causing increasing bit-flip rates in stored pages over the device lifetime.
Flash pages are typically 4 KB -- 16 KB; existing BCH or LDPC ECC controllers require dedicated out-of-band (OOB) storage that grows with error correction strength.

\textbf{MECC application:} MECC with explicit CRC metadata is stored in the OOB region of each flash page.
The $O(1/\sqrt{L})$ overhead scaling means that, for 16 KB pages ($L = 131,072$ bits), the same overhead budget that a BCH code uses to correct 50 errors can provide MECC correction of hundreds of errors.
The Feistel permutation makes MECC robust to the bursty error patterns that dominate aged flash.

\textbf{Key benefit:} Higher correction capacity per unit overhead than BCH for large page sizes; native burst-error resilience.

\subsection*{3. DRAM Reliability in Safety-Critical Systems}

\textbf{Context:} DRAM for automotive, avionics, and industrial control systems must meet strict reliability requirements (ISO 26262, DO-178C).
Single-bit ECC DRAM is standard but insufficient against multi-bit upsets from radiation or voltage transients.
Upgrading to multi-bit ECC doubles the DRAM footprint, which is often prohibitive.

\textbf{MECC application:} MECC with CRC-32 and large group sizes replaces the standard ECC DIMM with a scheme that corrects up to 200+ bit flips per 4096-bit block at roughly 50\% overhead --- substantially less than the 100\% overhead of a second DRAM chip.
A hardware accelerator implements the Feistel permutation and CRC hash computation in a memory controller, with the hash graph solver running in a dedicated microcode engine.

\textbf{Key benefit:} Multi-bit burst correction capability at lower hardware cost than full ECC DRAM.

\subsection*{4. Secure and Reliable Weight Storage for AI Accelerators}

\textbf{Context:} Custom AI accelerators (NPUs, TPUs, inference chips) integrate large on-chip weight SRAMs.
These dense SRAMs are vulnerable to soft errors under low-voltage operation needed for energy efficiency.
Any ECC overhead directly reduces the number of storable model parameters.

\textbf{MECC application:} The hybrid embedded/explicit MECC variant is applied layer-wise: sensitive layers (where weight distortion strongly affects accuracy) use explicit CRC storage; less sensitive layers use embedded storage.
The hybrid assignment minimizes accuracy impact while keeping total overhead below 5\%.
The MILP co-optimization runs offline during model compilation and adds no runtime cost.

\textbf{Key benefit:} Pareto-optimal tradeoff between storage overhead and model accuracy for each layer, with full burst-correction capability.

\subsection*{5. General-Purpose Configurable ECC for Storage Subsystems}

\textbf{Context:} Storage subsystem designers often require different protection levels for different data tiers (hot data vs.\ cold archive), but fixed-overhead BCH codes require a different controller for each protection level.

\textbf{MECC application:} A single MECC controller supports multiple protection profiles by software-configuring the CRC width and group size parameters.
No hardware changes are needed to switch between a lightweight 10\% overhead profile (CRC-8, large groups) and a high-protection 100\% overhead profile (CRC-32, small groups).
The same Feistel permutation engine and solver are reused across all profiles.

\textbf{Key benefit:} Single hardware implementation supporting a continuous range of overhead/protection tradeoffs, replacing multiple fixed-rate ECC controllers.
